\documentclass[tikz,border=6pt]{standalone}
\usepackage{ctex}
\usepackage{amsmath}
\usetikzlibrary{calc, arrows.meta}

\begin{document}
\begin{tikzpicture}[
  scale=1.0, thick,
]

% ─── 三栏 subplot：平行 / 重合 / 垂直 ───

% ── 栏 1：平行 ──
\begin{scope}[xshift=0cm]
  % 标题
  \node[font=\small\bfseries] at (1.4, 3.0) {平行 $l_1 \parallel l_2$};

  % l1
  \draw[black, thick] (0.0, 2.2) -- (2.8, 2.2);
  \node[font=\small] at (2.95, 2.2) {$l_1$};

  % l2
  \draw[black, thick] (0.0, 1.1) -- (2.8, 1.1);
  \node[font=\small] at (2.95, 1.1) {$l_2$};

  % 条件说明
  \node[font=\footnotesize, align=center] at (1.4, 0.35)
    {$k_1=k_2$，$b_1\neq b_2$};
\end{scope}

% ── 栏 2：重合 ──
\begin{scope}[xshift=4.2cm]
  \node[font=\small\bfseries] at (1.4, 3.0) {重合 $l_1 = l_2$};

  % l1（稍粗，蓝色偏移一点表示重叠）
  \draw[blue, dashed, thick] (0.0, 1.65) -- (2.8, 1.65);
  \draw[black, thick] (0.0, 1.65) -- (2.8, 1.65);
  \node[font=\small] at (3.05, 1.85) {$l_1$};
  \node[font=\small, blue] at (3.05, 1.45) {$l_2$};

  \node[font=\footnotesize, align=center] at (1.4, 0.35)
    {$k_1=k_2$，$b_1= b_2$};
\end{scope}

% ── 栏 3：垂直 ──
\begin{scope}[xshift=8.4cm]
  \node[font=\small\bfseries] at (1.4, 3.0) {垂直 $l_1 \perp l_2$};

  % l1：斜线
  \draw[black, thick] (0.3, 0.6) -- (2.5, 2.6);
  \node[font=\small] at (2.65, 2.65) {$l_1$};

  % l2：与 l1 垂直的斜线（斜率互为负倒数）
  \draw[red, thick] (0.3, 2.6) -- (2.5, 0.6);
  \node[font=\small, red] at (2.65, 0.55) {$l_2$};

  % 直角标记：交点大约在 (1.4, 1.6)
  \coordinate (P) at (1.4, 1.6);
  % 直角小方块：沿 l1 方向 (1,1)/sqrt(2) 和 l2 方向 (1,-1)/sqrt(2)
  \draw[gray, thin]
    ($(P) + 0.15*(0.707, 0.707)$) --
    ($(P) + 0.15*(0.707, 0.707) + 0.15*(0.707, -0.707)$) --
    ($(P) + 0.15*(0.707, -0.707)$);

  \node[font=\footnotesize, align=center] at (1.4, 0.35)
    {$k_1 k_2 = -1$};
\end{scope}

\end{tikzpicture}
\end{document}
